%-------------------------
% Cover Letter - Pure Storage Software Engineer (Core Engineering)
% Author: Ajay Venkatesh
%-------------------------

\documentclass[letterpaper,11pt]{article}

\usepackage[top=0.8in, bottom=0.8in, left=0.9in, right=0.9in]{geometry}
\usepackage[hidelinks]{hyperref}
\usepackage{lmodern}
\usepackage{parskip}
\usepackage{microtype}

\setlength{\parskip}{6pt}
\setlength{\parindent}{0pt}

\begin{document}

\begin{flushleft}
\textbf{\large Ajay Venkatesh} \\
Brooklyn, NY \\
+1 516 585 9013 \\
\href{mailto:av3855@nyu.edu}{av3855@nyu.edu}
\end{flushleft}

\vspace{10pt}

May 15, 2026

\vspace{6pt}

Hiring Manager \\
Pure Storage

\vspace{10pt}

Dear Hiring Manager,

My name is Ajay Venkatesh --- finishing my M.S. in Computer Engineering at NYU, with several years of production software experience: a few years at PwC in Bengaluru, a summer SDE internship at Amazon, and ongoing on-campus work at NYU's Novel AI Technologies. My background sits at the intersection of \textbf{distributed systems engineering, systems-level programming,} and \textbf{the patient end-to-end ownership} high-reliability infrastructure demands --- the same triangle the Core Engineering role at Pure Storage lives in.

\textbf{Python} is my daily driver for both backend services and ML pipelines --- I've shipped data preprocessing, model integration, and end-to-end ETL workflows across multiple projects. \textbf{C++} is where I cut my teeth in graduate coursework (Computer Architecture, System Design, Data Structures) and where I keep going back when an abstraction stops being honest about what's happening underneath. Clean, high-performance production code is the baseline I work to: at Amazon I shipped backend services in Java with low-level design patterns and unit/integration testing via \textbf{Mockito}, contributed under code-review discipline, and treated tests as part of the build.

On distributed systems specifically: at Amazon I built a high-throughput, event-driven processing system on \textbf{AWS ECS Fargate, SQS, SNS} with auto-scaling clusters, decoupling distributed message processing and materially reducing response latency. I provisioned resilient \textbf{Infrastructure-as-Code} on \textbf{AWS CDK} across multi-region CI/CD with rollback controls. At PwC I spent three years on a distributed backend on \textbf{Microsoft Azure} with event-driven processing and fault-tolerant retries. I'm fluent in \textbf{Docker} and \textbf{Kubernetes} for orchestration, and at home debugging production incidents through logs, metrics, and distributed traces.

What pulls me toward Pure Storage is the seriousness of the work. Six nines of uptime isn't a marketing line --- it's a constraint that shapes every architectural choice and every line of production code. Storage operating systems sit at a layer where correctness and performance both matter in real units, where bugs have real consequences for customers' data, and where engineering judgment compounds across decades of compatibility constraints. That's the kind of environment I want to grow inside of, alongside engineers who've been thinking about distributed consensus and data integrity at scale.

I'd welcome the chance to discuss further. Thanks for considering my application.

\vspace{12pt}

Sincerely, \\
Ajay Venkatesh

\end{document}
